% Unit 073 terjemahan Bahasa Indonesia dari chapter6.tex baris 321--556.
% Sumber: Methods of Algebra, Volume 2, commit
% 9a5803ff77dd3257484cb177f851a73770a59dd3 (CC BY 4.0).
% stable-unit-id: o014.aljabr2.chapter6.group-homology-cohomology
% Cakupan: Bagian 6.2 lengkap, ``Homologi dan Kohomologi Grup'';
% berhenti tepat sebelum Bagian 6.3 pada baris 557.

% segment-id: o014.li2.chapter6.group-homology-cohomology.h001
\section{Homologi dan Kohomologi Grup}\label{sec:group-coh-sub}
\hypertarget{o014.aljabr2.chapter6.group-homology-cohomology}{ }
% segment-id: o014.li2.chapter6.group-homology-cohomology.p001
Dalam bagian ini, gelanggang komutatif $\Bbbk$ dan grup $G$ tetap digunakan.

% segment-id: o014.li2.chapter6.group-homology-cohomology.d001
\begin{definisi}
	\index[sym1]{MGMG@$M^G$, $M_G$}
	Untuk setiap modul-$G$ $M$, definisikan modul-$\Bbbk$
% segment-id: o014.li2.chapter6.group-homology-cohomology.q001
	\begin{align*}
		M^G & := \left\{ m \in M : \forall g \in G, \; gm = m \right\}, \\
		M_G & := M \big/ \lrangle{gm - m : g \in G, \; m \in M}.
	\end{align*}
	Kita menyebut $M^G$ sebagai submodul \emph{invarian} dari $M$, dan $M_G$
	sebagai modul kuosien \emph{koinvarian} dari $M$.
\end{definisi}

% segment-id: o014.li2.chapter6.group-homology-cohomology.p002
Pengambilan invarian dan koinvarian masing-masing memberikan funktor
$\Bbbk$-linear $G\dcate{Mod}\to\Bbbk\dcate{Mod}$. Keduanya dapat dicirikan
sebagai adjoin bagi funktor inflasi $\mathrm{Infl}$ dari
Definisi~\ref{def:Res-Infl} dalam kasus $H=G$.

% segment-id: o014.li2.chapter6.group-homology-cohomology.pr001
\begin{proposisi}\label{prop:inv-coinv-infl}
	Tinjau funktor
	$\mathrm{Infl}:=\mathrm{Infl}^G_{\{1\}}:\Bbbk\dcate{Mod}\to
	G\dcate{Mod}$, yang memetakan setiap modul-$\Bbbk$ $N$ ke modul-$G$
	dengan aksi trivial $gy=y$. Terdapat dua pasangan adjoin
% segment-id: o014.li2.chapter6.group-homology-cohomology.g001
	\[\begin{tikzcd}[row sep=tiny]
		\mathrm{Infl}: \Bbbk\dcate{Mod} \arrow[shift left, r] & G\dcate{Mod} : (\cdot)^G , \arrow[shift left, l] \\
		(\cdot)_G: G\dcate{Mod} \arrow[shift left, r] & \Bbbk\dcate{Mod}: \mathrm{Infl}. \arrow[shift left, l]
	\end{tikzcd}\]
\end{proposisi}
% segment-id: o014.li2.chapter6.group-homology-cohomology.pf001
\begin{bukti}
	Misalkan $M$ modul-$G$ dan $N$ modul-$\Bbbk$. Menentukan homomorfisme
	modul-$G$ $\varphi:\mathrm{Infl}(N)\to M$ sama artinya dengan menentukan
	homomorfisme modul-$\Bbbk$ $\varphi:N\to M$ sedemikian sehingga
	$\varphi(y)=g\varphi(y)$ selalu berlaku; ini juga sama artinya dengan
	menentukan homomorfisme modul-$\Bbbk$ $\varphi:N\to M^G$.
	
	Di pihak lain, menentukan homomorfisme modul-$\Bbbk$ $\psi:M_G\to N$
	sama artinya dengan menentukan homomorfisme modul-$\Bbbk$
	$\varphi:M\to N$ sedemikian sehingga $\varphi(gx-x)=0$ selalu berlaku.
	Syarat terakhir ekuivalen dengan $\varphi(gx)=g\varphi(x)$, dengan
	aksi $G$ pada ruas kanan didefinisikan oleh struktur modul-$G$
	$\mathrm{Infl}(N)$.
\end{bukti}

% segment-id: o014.li2.chapter6.group-homology-cohomology.p003
Funktor $(\cdot)^G$ dan $(\cdot)_G$ juga dapat dipahami dari sudut pandang
modul-$\Bbbk[G]$. Berikut ini, melalui homomorfisme augmentasi dalam
Catatan~\ref{rem:augmentation-kG}, kita memandang $\Bbbk$ sebagai
bimodul-$(\Bbbk[G],\Bbbk[G])$; dengan kata lain, aksi $G$ pada kedua sisi
$\Bbbk$ bersifat trivial.

% segment-id: o014.li2.chapter6.group-homology-cohomology.pr002
\begin{proposisi}\label{prop:inv-coinv-Hom}
	Dengan mengidentifikasi modul-$G$ dan modul kiri-$\Bbbk[G]$, terdapat
	isomorfisme alami antarfunktor
% segment-id: o014.li2.chapter6.group-homology-cohomology.q002
	\begin{gather*}
		(\cdot)^G \simeq \Hom_G(\Bbbk, \cdot), \quad  (\cdot)_G \simeq \Bbbk \dotimes{\Bbbk[G]} (\cdot).
	\end{gather*}
\end{proposisi}
% segment-id: o014.li2.chapter6.group-homology-cohomology.pf002
\begin{bukti}
	Tinjau modul-$G$ $M$. Berdasarkan Proposisi~\ref{prop:inv-coinv-infl},
	menentukan homomorfisme modul-$G$ $\varphi:\Bbbk\to M$ sama artinya
	dengan menentukan homomorfisme modul-$\Bbbk$ $\varphi:\Bbbk\to M^G$.
	Yang terakhir sama artinya dengan menentukan unsur $M^G$, yaitu
	$\varphi(1)$. Ini memberikan isomorfisme pertama.
	
	Untuk isomorfisme kedua, identifikasikan terlebih dahulu $\Bbbk$ dengan
	$\Bbbk[G]/\mathfrak{I}$, dengan ideal augmentasi
	$\mathfrak{I}=\bigoplus_{g\neq1_G}\Bbbk(g-1_G)$. Maka
	$\Bbbk\dotimes{\Bbbk[G]}M\simeq M/\mathfrak{I}M=M_G$, melalui pemetaan
	$t\otimes x$ ke citra $tx\in M$ dalam $M_G$. Terbukti.
\end{bukti}

% segment-id: o014.li2.chapter6.group-homology-cohomology.p004
Dengan demikian, adjungsi dalam Proposisi~\ref{prop:inv-coinv-infl} dapat
dipahami sebagai kasus khusus Proposisi~\ref{prop:pullback-two-adjoint}, yaitu
dengan mengambil homomorfisme trivial $\varphi:G\twoheadrightarrow\{1\}$.

% segment-id: o014.li2.chapter6.group-homology-cohomology.co001
\begin{korolari}\label{prop:Ind-invariant}
	Untuk setiap homomorfisme grup $\varphi:H\to G$, terdapat isomorfisme funktor
% segment-id: o014.li2.chapter6.group-homology-cohomology.q003
	\[\begin{array}{rlrl}
		\Hom_H(\Bbbk[G], \cdot)^G & \rightiso (\cdot)^H, & (\cdot)_H & \rightiso (\Bbbk[G] \dotimes{\Bbbk[H]} \cdot)_G , \\
		\varphi & \mapsto \varphi(1_G), & (x \;\text{citranya}) & \mapsto (1_G \otimes x \;\text{citranya}) .
	\end{array}\]
	Kedua ruas merupakan funktor dari $H\dcate{Mod}$ ke
	$\Bbbk\dcate{Mod}$. Ruas kiri melibatkan funktor
	$\Hom_H(\Bbbk[G],\cdot):H\dcate{Mod}\to G\dcate{Mod}$ yang diperkenalkan
	dalam Proposisi~\ref{prop:pullback-two-adjoint}.
\end{korolari}
% segment-id: o014.li2.chapter6.group-homology-cohomology.pf003
\begin{bukti}
	Berikut ini, misalkan $N$ modul-$H$. Isomorfisme pertama mengikuti dari
	adjungsi dalam Proposisi~\ref{prop:pullback-two-adjoint}:
% segment-id: o014.li2.chapter6.group-homology-cohomology.q004
	\[ \Hom_G(\Bbbk, \Hom_H(\Bbbk[G], N)) \simeq \Hom_H(\varphi^* \Bbbk, N) = \Hom_H(\Bbbk, N). \]
	
	Isomorfisme kedua mengikuti dari kendala asosiativitas hasil kali tensor
	$\Bbbk\dotimes{\Bbbk[G]}(\Bbbk[G]\dotimes{\Bbbk[H]}N)
	\simeq\Bbbk\dotimes{\Bbbk[H]}N$. Verifikasi pemetaan eksplisitnya
	diserahkan kepada pembaca. Sebagai alternatif, pernyataan tersebut juga
	dapat dibuktikan langsung dengan menunjukkan bahwa pemetaan yang
	ditampilkan terdefinisi dengan baik dan dapat dibalik.
\end{bukti}

% segment-id: o014.li2.chapter6.group-homology-cohomology.p005
Pasangan-pasangan adjoin dalam Proposisi~\ref{prop:inv-coinv-infl} juga
memastikan bahwa $(\cdot)_G$ eksak kanan dan $(\cdot)^G$ eksak kiri; hal ini
juga mudah diverifikasi secara langsung.

% segment-id: o014.li2.chapter6.group-homology-cohomology.d002
\begin{definisi}[Homologi dan kohomologi grup]\label{def:group-coh-ho}
	\index{homologi grup dan kohomologi grup (homology/cohomology of groups)}
	\index[sym1]{HnGM@$\Hm_n(G, M)$, $\Hm^n(G, M)$}
	Misalkan $G$ suatu grup. \emph{Homologi grup} dengan koefisien dalam $\Bbbk$
	didefinisikan sebagai funktor turunan kiri dari $(\cdot)_G$:
% segment-id: o014.li2.chapter6.group-homology-cohomology.q005
	\[ \Hm_n(G, \cdot) := \mathrm{L}_n (\cdot)_G: G\dcate{Mod} \to \Bbbk\dcate{Mod}, \]
	sedangkan \emph{kohomologi grup} didefinisikan sebagai funktor turunan
	kanan dari $(\cdot)^G$:
% segment-id: o014.li2.chapter6.group-homology-cohomology.q006
	\[ \Hm^n(G, \cdot) := \mathrm{R}^n (\cdot)^G: G\dcate{Mod} \to \Bbbk\dcate{Mod}, \]
	dengan $n\in\Z$. Keduanya hanya tak nol untuk $n\geq0$. Definisi serupa
	juga berlaku bagi modul kanan-$G$.
\end{definisi}

% segment-id: o014.li2.chapter6.group-homology-cohomology.p006
Mengikuti kebiasaan dalam banyak literatur, jika koefisien tidak disebutkan,
diandaikan $\Bbbk=\Z$. Dalam hal ini, modul-$G$ sama artinya dengan grup
aditif yang dilengkapi aksi kiri $G$, dengan aksi grup yang mempertahankan
penjumlahan.

% segment-id: o014.li2.chapter6.group-homology-cohomology.e001
\begin{contoh}\label{eg:group-H1}
	Tinjau ideal augmentasi
	$\mathfrak{I}=\bigoplus_{g\neq1_G}\Bbbk(g-1_G)$ dari $\Bbbk[G]$. Untuk
	setiap modul-$G$ $M$, terdapat
	$\Hm_0(G,M)\simeq M_G\simeq M/\mathfrak{I}M$; khususnya,
	$\Hm_0(G,\mathfrak{I})\simeq\mathfrak{I}/\mathfrak{I}^2$.

	Sekarang tinjau barisan eksak pendek
	$0\to\mathfrak{I}\to\Bbbk[G]\to\Bbbk\to0$, yang morfisme terakhirnya
	adalah homomorfisme augmentasi. Karena $\Bbbk[G]$ merupakan modul
	projektif, teknik penggeseran dimensi
	(Proposisi~\ref{prop:dimension-shifting}) menunjukkan
% segment-id: o014.li2.chapter6.group-homology-cohomology.q007
	\[ \Hm_n(G, \Bbbk) \simeq \begin{cases}
		\Hm_{n-1}(G, \mathfrak{I}), & n \geq 2 \\
		\Ker\left[ \mathfrak{I}/\mathfrak{I}^2 \to \Bbbk[G]/\mathfrak{I} \right] = \mathfrak{I}/\mathfrak{I}^2 , & n = 1.
	\end{cases}\]
\end{contoh}

% segment-id: o014.li2.chapter6.group-homology-cohomology.p007
Berdasarkan isomorfisme antara $\Bbbk[G]\dcate{Mod}$ dan $G\dcate{Mod}$,
serta Proposisi~\ref{prop:inv-coinv-Hom}, homologi dan kohomologi grup juga
dapat ditafsirkan melalui funktor $\Tor$ dan $\Ext$ yang diperkenalkan dalam
\S\ref{sec:Ext-Tor} sebagai berikut:
% segment-id: o014.li2.chapter6.group-homology-cohomology.q008
\begin{equation}\label{eqn:grp-Hm-Tor-Ext}\begin{aligned}
	\Hm_n(G, M) & \simeq \Tor^G_n(\Bbbk, M) := \Tor^{\Bbbk[G]}_n(\Bbbk, M), \\
	\Hm^n(G, M) & \simeq \Ext_G^n(\Bbbk, M) := \Ext_{\Bbbk[G]}^n(\Bbbk, M).
\end{aligned}\end{equation}

% segment-id: o014.li2.chapter6.group-homology-cohomology.p008
Penafsiran ini kadang-kadang mereduksi persoalan terkait menjadi hasil yang
telah diketahui dalam teori modul. Berikut ini satu contoh mendasar yang
tidak sepenuhnya trivial dan melibatkan barisan spektral yang menghubungkan
$\Ext$ dan $\Tor$.

% segment-id: o014.li2.chapter6.group-homology-cohomology.pr003
\begin{proposisi}[Teorema koefisien universal untuk homologi dan kohomologi grup]\label{prop:group-UCT}
	\index{teorema koefisien universal}
	Ambil $\Bbbk=\Z$. Lengkapi modul-$\Z$ $M$ dengan aksi trivial $G$ sehingga
	menjadi modul-$G$. Untuk setiap $n\in\Z_{\geq0}$, terdapat barisan eksak pendek
	kanonik
% segment-id: o014.li2.chapter6.group-homology-cohomology.g002
	\[\begin{tikzcd}[row sep=tiny, column sep=small]
		0 \arrow[r] & \Ext^1_{\Z}\left( \Hm_{n-1}(G, \Z), M \right) \arrow[r] & \Hm^n(G, M) \arrow[r] & \Hom_{\Z}\left( \Hm_n(G, \Z), M \right) \arrow[r] & 0, \\
		0 \arrow[r] & M \dotimes{\Z} \Hm_n(G, \Z) \arrow[r] & \Hm_n(G, M) \arrow[r] & \Tor^{\Z}_1(M, \Hm_{n-1}(G, \Z)) \arrow[r] & 0;
	\end{tikzcd}\]
	Untuk $n=0$, digunakan konvensi $\Hm_{-1}(G,\Z)=0$.
\end{proposisi}
% segment-id: o014.li2.chapter6.group-homology-cohomology.pf004
\begin{bukti}
	Kita boleh mengandaikan $n\geq1$. Pertama, tinjau kasus kohomologi. Setelah
	pembaca menguasai kompleks kokrantai standar dan kompleks rantai yang akan
	diperkenalkan kemudian, barisan eksak pendek tersebut mudah diperoleh
	dengan memasukkan kompleks rantai
	$C:=\mathsf{L}\dotimes{\Bbbk[G]}\Z$ dan modul-$\Z$ $M$ ke teorema
	koefisien universal untuk kohomologi yang diperkenalkan dalam latihan
	\CHref{sec:cplx}. Perinciannya diserahkan kepada pembaca yang berminat.
	
	Di sini kita sajikan argumen lain. Dalam
	Contoh~\ref{eg:change-of-rings-SS-bis}, ambil homomorfisme gelanggang
	$R:=\Z[G]\to S:=\Z$ sebagai augmentasi; selain itu, ambil modul kiri-$R$
	$Y=\Z$ (yaitu modul-$G$ trivial) dan modul kiri-$S$ $X=M$. Barisan eksak
	pendek dalam contoh tersebut lalu menjadi bentuk yang dicari.
	
	Untuk kasus homologi, dengan cara yang sama kita dapat memasukkan kompleks
	rantai standar ke teorema Künneth untuk homologi
	\ref{prop:Kunneth-homology}, dengan perincian diserahkan kepada pembaca,
	atau memahaminya melalui Contoh~\ref{eg:change-of-rings-SS-bis}.
\end{bukti}

% segment-id: o014.li2.chapter6.group-homology-cohomology.p009
Mengikuti konstruksi yang telah diperkenalkan dalam
\S\ref{sec:derived-primer}, kita juga dapat mengambil funktor turunan terkait
untuk kompleks rantai yang terbatas di bawah (atau kompleks kokrantai yang
terbatas di bawah) $M$. Funktor tersebut disebut hiperhomologi grup (atau
hiperkohomologi grup). Lebih lanjut, dalam kategori turunan dapat ditinjau
% segment-id: o014.li2.chapter6.group-homology-cohomology.q009
\[ \Bbbk \otimesL[{\Bbbk[G]}] M, \quad \RHom_G(\Bbbk, M) := \RHom_{\Bbbk[G]}(\Bbbk, M). \]

% segment-id: o014.li2.chapter6.group-homology-cohomology.p010
Berdasarkan \eqref{eqn:grp-Hm-Tor-Ext} dan fakta bahwa funktor $\Ext$ (atau
$\Tor$) dapat dihitung dengan mengambil resolusi projektif pada variabel
pertama, resolusi bebas $\mathsf{L}\to\Bbbk$ dari
Proposisi~\ref{prop:trivial-mod-resolution} langsung memberikan kompleks
(atau kompleks rantai) konkret untuk menghitung kohomologi (atau homologi)
grup. Definisi~\ref{def:resolution-trivial-module} menyediakan dua cara untuk
menjadikan $\mathsf{L}$ kompleks modul kiri-$G$ atau modul kanan-$G$: gunakan
$\partial_n:\mathsf{L}_n\to\mathsf{L}_{n-1}$ atau gunakan
	$\partial'_n:\mathsf{L}_n\to\mathsf{L}_{n-1}$. Buku ini akan memilih salah
	satunya sesuai konteks agar diperoleh bentuk kompleks yang standar,
tetapi pilihan ini tidak memengaruhi kohomologi kompleks tersebut.

% segment-id: o014.li2.chapter6.group-homology-cohomology.pr004
\begin{proposisi}[Menghitung kohomologi grup dengan kompleks kokrantai standar]\label{prop:groupcoh-cochain}
	\index{kompleks kokrantai standar (standard complex)}
	\index[sym1]{CGM@$C(G, M)$}
	Untuk setiap modul-$G$ $M$, definisikan \emph{kompleks kokrantai standar}
	$C(G,M)=(C^n(G,M))_n$, yang terdiri atas modul-$\Bbbk$, sebagai berikut:
% segment-id: o014.li2.chapter6.group-homology-cohomology.q010
	\begin{gather*}
		C^n(G, M) := \left\{ \text{pemetaan}\; f: G^n \to M \right\}, \quad n \in \Z_{\geq 0},
	\end{gather*}
	Setiap suku menjadi modul-$\Bbbk$ melalui penjumlahan dan perkalian skalar
	titik demi titik; suku-suku berderajat negatif didefinisikan sebagai $0$;
	dan $d^n:C^n(G,M)\to C^{n+1}(G,M)$ didefinisikan oleh
% segment-id: o014.li2.chapter6.group-homology-cohomology.q011
	\begin{multline*}
		(d^n f)(g_1, \ldots, g_{n+1}) = g_1 f(g_2, \ldots, g_{n+1}) \\
		+ \sum_{k=1}^n (-1)^k f(\ldots, g_k g_{k+1}, \ldots) + (-1)^{n+1} f(g_1, \ldots, g_n).
	\end{multline*}

	Untuk setiap $n$, terdapat isomorfisme kanonik
	$\Hm^n(C(G,M))\simeq\Hm^n(G,M)$. Lebih tepatnya, $C(G,M)$ mewakili
	$\RHom_G(\Bbbk,M)$ dalam kategori turunan.
\end{proposisi}
% segment-id: o014.li2.chapter6.group-homology-cohomology.pf005
\begin{bukti}
	Tinjau kompleks $\mathsf{L}=(\mathsf{L}_n,\partial_n)_n$. Berdasarkan
	resolusi bebas modul kiri-$G$ $\mathsf{L}\to\Bbbk$,
	$\RHom_G(\Bbbk,M)$ diwakili oleh kompleks $\Hom$ berikut
	(Definisi~\ref{def:Hom-cplx}):
% segment-id: o014.li2.chapter6.group-homology-cohomology.q012
	\begin{equation*}
		\Hom^\bullet(\mathsf{L}, M) = \left[ \cdots \to \Hom_{\Bbbk[G]}(\mathsf{L}_n, M) \xrightarrow{(-1)^{n+1} \partial_{n+1}^*} \Hom_{\Bbbk[G]}(\mathsf{L}_{n+1}, M) \to \cdots \right] ,
	\end{equation*}
	yang derajatnya adalah $\ldots,-n,-n-1,\ldots$; tanda $(-1)^{n+1}$
	berasal dari definisi umum kompleks $\Hom$.
	
	Untuk setiap $n\geq0$, gunakan basis
	$(g_1|\cdots|g_n)$ bagi modul kiri-$\Bbbk[G]$ bebas $\mathsf{L}_n$
	seperti dalam \eqref{eqn:L-basis}. Dengan demikian,
% segment-id: o014.li2.chapter6.group-homology-cohomology.g003
	\[\begin{tikzcd}[row sep=tiny, column sep=tiny]
		\Hom_G(\mathsf{L}_n, M) \arrow[r, "\sim"] & C^n(G, M) \\
		\phi \arrow[mapsto, r] & {\left[ f: (g_1, \ldots, g_n) \mapsto \phi\left( (g_1 | \cdots | g_n) \right)\right]} ,
	\end{tikzcd}\]
	yang jelas merupakan isomorfisme modul-$\Bbbk$. Langkah kuncinya ialah
	mengidentifikasi $(-1)^{n+1}\partial_{n+1}^*$ pada $C^n(G,M)$. Tinjau
	$(-1)^{n+1}\partial_{n+1}^*\phi=(-1)^{n+1}\phi\circ\partial_{n+1}$.
	Berdasarkan \eqref{eqn:L-basis-partial}, pemetaan ini mengirim
	$(g_1|\cdots|g_{n+1})\in\mathsf{L}_{n+1}$ ke
% segment-id: o014.li2.chapter6.group-homology-cohomology.q013
	\begin{multline*}
		\phi\left( g_1 (g_2 | \cdots | g_{n+1})\right) + \sum_{k=1}^n (-1)^k \phi\left((\cdots | g_k g_{k+1}| \cdots)\right) + \cdots + (-1)^{n+1} \phi\left((g_1 | \cdots |g_n)\right) \\
		= g_1 f(g_2, \ldots, g_{n+1}) + \sum_{k=1}^n (-1)^k f(\ldots, g_k g_{k+1}, \ldots) + (-1)^{n+1} f(g_1, \ldots, g_n).
	\end{multline*}
	Inilah tepatnya $(d^nf)(g_1,\ldots,g_{n+1})$ dalam
	pernyataan.\footnote{Koreksi penerjemah: sumber menuliskan argumen hingga
	$g_n$, tetapi $d^nf\in C^{n+1}(G,M)$ dan perhitungan di atas memakai
	$(g_1,\ldots,g_{n+1})$.}
\end{bukti}

% segment-id: o014.li2.chapter6.group-homology-cohomology.p011
Kohomologi modul kanan-$G$ $M$ mempunyai deskripsi serupa, yang tidak kita
ulangi. Berikutnya, kita nyatakan hasil yang bersesuaian bagi homologi grup
dari modul kiri-$G$; buktinya langsung.

% segment-id: o014.li2.chapter6.group-homology-cohomology.pr005
\begin{proposisi}[Menghitung homologi grup dengan kompleks rantai standar]\label{prop:grouph-chain}
	Untuk setiap modul-$G$ $M$ dan $n\in\Z_{\geq0}$, terdapat isomorfisme
	kanonik
% segment-id: o014.li2.chapter6.group-homology-cohomology.q014
	\[ \Hm_n(G, M) \simeq \Hm_n\left( \mathsf{L} \dotimes{\Bbbk[G]} M\right), \quad n \in \Z_{\geq 0}, \]
	dengan $(\mathsf{L}_n,\partial'_n)_n$ digunakan sebagai struktur kompleks
	rantai modul kanan-$G$. Ambil basis
	$(g_1|\cdots|g_n)^\dagger$ bagi modul kanan-$\Bbbk[G]$ bebas
	$\mathsf{L}_n$ seperti dalam \eqref{eqn:L-basis-right}. Maka
	$\mathsf{L}_n\dotimes{\Bbbk[G]}M$ diidentifikasi dengan jumlah langsung
	\footnote{Koreksi penerjemah: sumber menuliskan $\mathsf{L}^n$, sedangkan
	kompleks rantai dan basis pada kalimat yang sama menggunakan $\mathsf{L}_n$.}
	$M^{\oplus G^n}$, dan unsur pada suku jumlah langsung yang bersesuaian
	dengan $(g_1,\ldots,g_n)\in G^n$ dapat dituliskan secara tunggal sebagai
% segment-id: o014.li2.chapter6.group-homology-cohomology.q015
	\[ (g_1 | \cdots | g_n)^\dagger \otimes x, \quad x \in M, \]
	sedangkan, berdasarkan \eqref{eqn:L-basis-partial-right},
	$\partial'_n\otimes\identity_M:\mathsf{L}_n\dotimes{\Bbbk[G]}M\to
	\mathsf{L}_{n-1}\dotimes{\Bbbk[G]}M$ dinyatakan oleh
% segment-id: o014.li2.chapter6.group-homology-cohomology.q016
	\begin{multline*}
		(g_1 | \cdots | g_n)^\dagger \otimes x \mapsto \\
		(g_2 | \cdots | g_n)^\dagger \otimes x + \sum_{k=1}^{n-1} (-1)^k (\cdots | g_k g_{k+1} | \cdots)^\dagger \otimes x \\
		+ (-1)^n (g_1 | \cdots | g_{n-1})^\dagger \otimes g_n x.
	\end{multline*}

	Untuk membandingkannya dengan Proposisi~\ref{prop:groupcoh-cochain}, kita
	menuliskan kompleks rantai di atas sebagai $(C_n(G,M))_n$ dan menyebutnya
	\emph{kompleks rantai standar}. Dalam kategori turunan, kompleks ini
	memberikan $\Bbbk\otimesL[{\Bbbk[G]}]M$.
\end{proposisi}

% segment-id: o014.li2.chapter6.group-homology-cohomology.p012
Contoh~\ref{eg:group-H-comonad} nanti akan menjelaskan hubungan antara
kompleks rantai standar dan konstruksi bar.

% segment-id: o014.li2.chapter6.group-homology-cohomology.p013
Tentu terdapat pula deskripsi yang bersesuaian bagi homologi modul kanan-$G$
$M$. Tinjau $M\dotimes{\Bbbk[G]}\mathsf{L}$. Terdapat
$M\dotimes{\Bbbk[G]}\mathsf{L}_n\simeq M^{\oplus G^n}$, dan unsur pada suku
jumlah langsung yang bersesuaian dengan $(g_1,\ldots,g_n)$ dapat dituliskan
secara tunggal sebagai
% segment-id: o014.li2.chapter6.group-homology-cohomology.q017
\[ x \otimes (g_1 | \cdots | g_n), \quad x \in M, \]
sedangkan pemetaan $\identity_M\otimes\partial'_n$ adalah
% segment-id: o014.li2.chapter6.group-homology-cohomology.q018
\begin{multline*}
	x \otimes (g_1 | \cdots | g_n) \mapsto \\
	x g_1 (g_2 | \cdots | g_n) + \sum_{k=1}^{n-1} (-1)^k x \otimes (\cdots | g_k g_{k+1} | \cdots) \\
	+ (-1)^n x \otimes (g_1 | \cdots | g_{n-1}).
\end{multline*}
Kompleks rantai standar yang bersesuaian tetap dituliskan $(C_n(G,M))_n$.

% segment-id: o014.li2.chapter6.group-homology-cohomology.p014
Pembaca yang mengenal isi \S\ref{sec:HH} akan mudah mengenali
kompleks-kompleks ini sebagai kasus khusus kompleks kokrantai dan kompleks
rantai Hochschild. Berikut pernyataan terperincinya.

% segment-id: o014.li2.chapter6.group-homology-cohomology.pr006
\begin{proposisi}[Hubungan dengan teori Hochschild]
	Misalkan $R:=\Bbbk[G]$. Untuk modul kiri-$G$ (atau modul kanan-$G$) $M$,
	tuliskan $M_+$ (atau ${}_+M$) untuk bimodul-$(R,R)$ yang ditentukan oleh
% segment-id: o014.li2.chapter6.group-homology-cohomology.q019
	\[ g_1 x g_2 := g_1 x \; (\text{atau}\; g_1 x g_2 = x g_2 ), \quad g_1 , g_2 \in G, \; x \in M. \]
	Terdapat isomorfisme kanonik yang kompatibel dengan barisan eksak panjang:
% segment-id: o014.li2.chapter6.group-homology-cohomology.q020
	\[ \Hm^n(G, M) \simeq \HHm^n(R, M_+), \quad (\text{atau}\; \Hm_n(G, M) \simeq \HHm_n(R, {}_+ M)), \]
	dan bahkan kompleks $(C^n(G,M))_n$ (atau kompleks rantai
	$(C_n(G,M))_n$) isomorfik dengan $(C^n(R,M_+))_n$ (atau
	$(C_n(R,{}_+M))_n$) dari \eqref{eqn:HH-cplx}.
\end{proposisi}
% segment-id: o014.li2.chapter6.group-homology-cohomology.pf006
\begin{bukti}
	Perhatikan bahwa $R^{\otimes n}$ merupakan modul-$\Bbbk$ bebas dengan
	basis $G^n$: unsur $(g_1,\ldots,g_n)$ bersesuaian dengan
	$g_1\otimes\cdots\otimes g_n\in R^{\otimes n}$. Tinggal membandingkan
	kompleks kokrantai standar dengan kompleks dari \eqref{eqn:HH-cplx}.
	Secara konkret, dalam kasus kohomologi, kita memetakan $f:G^n\to M$ ke
	pemetaan $\Bbbk$-linear $n$-lipat $R^n\to M$ yang bersesuaian. Dalam
	kasus homologi, kita memetakan
	$x\otimes(g_1|\cdots|g_n)^\dagger$ ke
	$(x|g_1|\cdots|g_n)\in C_n(R,{}_+M)$. Perbandingan selanjutnya diserahkan
	kepada pembaca.
\end{bukti}

% segment-id: o014.li2.chapter6.group-homology-cohomology.p015
Kembali ke pokok bahasan, selanjutnya modul-$G$ yang dibahas selalu diandaikan
sebagai modul kiri.

% segment-id: o014.li2.chapter6.group-homology-cohomology.e002
\begin{contoh}\label{eg:group-H1-further}
	Berdasarkan pembahasan di atas, tinjau
% segment-id: o014.li2.chapter6.group-homology-cohomology.q021
	\[ \Hm_1(G, \Bbbk) \simeq \frac{\Ker\left[ C_1(G, \Bbbk) \to C_0(G, \Bbbk) \right]}{\Image\left[ C_2(G, \Bbbk) \to C_1(G, \Bbbk) \right]}. \]
	
	Setiap unsur $C_1(G,\Bbbk)$ dapat dituliskan secara tunggal sebagai
	kombinasi $\Bbbk$-linear dari unsur-unsur
	$[g]:=(g)^\dagger\otimes1$ ($g\in G$). Perhatikan bahwa
	$C_0(G,\Bbbk)\simeq\Bbbk$ dan
% segment-id: o014.li2.chapter6.group-homology-cohomology.q022
	\[ \left( \partial'_1 \otimes \identity_{\Bbbk}\right)[g] = 1 - 1 = 0, \quad \left( \partial'_2 \otimes \identity_{\Bbbk}\right)((g_1 | g_2)^\dagger \otimes 1) = [g_2] - [g_1 g_2] + [g_1]. \]
	Karena itu, $\Hm_1(G,\Bbbk)$ isomorfik dengan modul-$\Bbbk$ kuosien
% segment-id: o014.li2.chapter6.group-homology-cohomology.q023
	\[ Q := \bigoplus_{g \in G} \Bbbk \cdot [g] \bigg/ \lrangle{ [g_1 g_2] - [g_1] + [g_2] : g_1, g_2 \in G }; \]
	Pemetaan $g\mapsto[g]$ menginduksi homomorfisme grup $a:G\to(Q,+)$.
	Pemetaan ini mempunyai sifat universal berikut: untuk setiap modul-$\Bbbk$
	$N$ dan homomorfisme grup $b:G\to(N,+)$, terdapat tepat satu homomorfisme
	modul-$\Bbbk$ $\varphi:Q\to N$ sedemikian sehingga $b=\varphi a$.
	Pembaca dipersilakan segera memverifikasi sifat universal ini dan
	menurunkan darinya isomorfisme kanonik modul-$\Bbbk$
% segment-id: o014.li2.chapter6.group-homology-cohomology.g004
	\[\begin{tikzcd}[row sep=tiny]
		G_{\mathrm{ab}} \dotimes{\Z} \Bbbk \arrow[r, "\sim"] & \Hm_1(G, \Bbbk) \\
		(g \bmod G_{\mathrm{der}}) \otimes 1 \arrow[mapsto, r] & {[g] \;\text{kelasnya}},
	\end{tikzcd}\]
	dengan $G_{\mathrm{ab}}:=G/G_{\mathrm{der}}$ abelianisasi $G$, dan
	$G_{\mathrm{der}}$ subgrup turunan $G$, seperti dalam
	\cite[Definisi 4.7.2]{Li1}. Pembaca yang mengetahui latar topologis
	homologi grup akan mengenali kaitannya dengan teorema Hurewicz dalam teori
	homologi.
	
	Secara alami, kita juga ingin memahami komposisi isomorfisme
	$\Hm_1(G,\Bbbk)\rightiso\mathfrak{I}/\mathfrak{I}^2$ dari
	Contoh~\ref{eg:group-H1} dengan isomorfisme di atas. Kita menyatakan bahwa
	komposisi tersebut diberikan secara konkret oleh
% segment-id: o014.li2.chapter6.group-homology-cohomology.g005
	\[\begin{tikzcd}[row sep=tiny]
		G_{\mathrm{ab}} \otimes \Bbbk \arrow[r, "\sim"] & \mathfrak{I}/\mathfrak{I}^2 \\
		(g \bmod G_{\mathrm{der}}) \otimes 1 \arrow[mapsto, r] & 1_G - g \;\bmod \mathfrak{I}^2.
	\end{tikzcd}\]
	
	Mengapa demikian? Isomorfisme dalam Contoh~\ref{eg:group-H1} berasal dari
	$\Hm_1(G,\Bbbk)\to\Hm_0(G,\mathfrak{I})\simeq
	\mathfrak{I}/\mathfrak{I}^2$ yang diinduksi oleh barisan eksak pendek
	$0\to\mathfrak{I}\to\Bbbk[G]\to\Bbbk\to0$. Untuk itu, tuliskan diagram
	komutatif dengan baris-baris eksak berikut:
% segment-id: o014.li2.chapter6.group-homology-cohomology.g006
	\[\begin{tikzcd}
		& & {(g)^\dagger \otimes 1_G} \arrow[phantom, d, "\in" description, sloped] \arrow[mapsto, r] & {[g] = (g)^\dagger \otimes 1} \arrow[phantom, d, "\in" description, sloped] & \\
		0 \arrow[r] & C_1(G, \mathfrak{I}) \arrow[r] \arrow[d] & C_1(G, \Bbbk[G]) \arrow[r] \arrow[d, "{\partial'_1 \otimes \identity_{\Bbbk[G]}}"] & C_1(G, \Bbbk) \arrow[r] \arrow[d, "{\partial'_0 \otimes \identity_{\Bbbk[G]} = 0}"] & 0 \\
		0 \arrow[r] & C_0(G, \mathfrak{I}) \arrow[r] & C_0(G, \Bbbk[G]) \arrow[r] & C_0(G, \Bbbk) \arrow[r] & 0 ,
	\end{tikzcd}\]
	Pemetaan $\partial'_1\otimes\identity_{\Bbbk[G]}$ mengirim
	$(g)^\dagger\otimes1_{\Bbbk[G]}$ ke
	$1_G-g\in C_0(G,\mathfrak{I})=\mathfrak{I}$, dan dengan demikian
	menentukan unsur $\mathfrak{I}/\mathfrak{I}^2\simeq
	\Hm_0(G,\mathfrak{I})$.\footnote{Koreksi penerjemah: sumber memakai
	$C^0$ dan $\Hm^0$, tetapi diagram dan morfisme penghubung yang sedang
	dihitung adalah homologis, sehingga indeks yang benar ialah $C_0$ dan
	$\Hm_0$.} Inilah tepatnya konstruksi dalam lema ular.
\end{contoh}

% segment-id: o014.li2.chapter6.group-homology-cohomology.r001
\begin{catatan}\label{rem:Gmod-k}
	Bentuk kompleks kokrantai standar (atau kompleks rantai standar)
	menunjukkan bahwa, sebagai grup aditif, $\Hm^n(G,M)$ (atau
	$\Hm_n(G,M)$) sepenuhnya ditentukan oleh struktur aditif $M$; perkalian
	dari $\Bbbk$ hanya memberinya struktur modul-$\Bbbk$. Hal ini menunjukkan
	bahwa kasus $\Bbbk=\Z$ telah menangkap esensi kohomologi (atau homologi)
	grup, meskipun mengizinkan koefisien umum tidak menimbulkan kesulitan
	tambahan.
\end{catatan}

% segment-id: o014.li2.chapter6.group-homology-cohomology.r002
\begin{catatan}[Versi ternormalisasi]\label{rem:nor-cochain}
	\index{kompleks kokrantai standar!ternormalisasi (normalized)}
	\index[sym1]{CGMbar@$\overline{C}(G, M)$}
	Untuk menghitung $\Hm^n(G,M)$ dan $\Hm_n(G,M)$, kita juga dapat memakai
	resolusi bebas ternormalisasi
	$\overline{\mathsf{L}}\to\Bbbk$ dari
	Proposisi~\ref{prop:trivial-mod-nor-resolution}. Kompleks kokrantai dan
	kompleks rantai yang bersesuaian masing-masing dituliskan
	$(\overline{C}^n(G,M))_n$ dan $(\overline{C}_n(G,M))_n$, serta disebut
	kompleks kokrantai standar ternormalisasi dan kompleks rantai standar
	ternormalisasi. Keduanya berturut-turut merupakan subkompleks dan kompleks
	kuosien dari versi semula. Sebagai contoh, untuk kasus kohomologi,
	deskripsi $\mathsf{L}_n$ dalam \eqref{eqn:nor-cochain} memberikan
% segment-id: o014.li2.chapter6.group-homology-cohomology.q024
	\[ \overline{C}^n(G, M) = \left\{\begin{array}{r|l}
		f \in C^n(G, M) & \exists 1 \leq k \leq n, \; g_k = 1_G \\
		& \implies f(g_1, \ldots, g_n) = 0
	\end{array}\right\}. \]
\end{catatan}
